
\subsection{Project Components}
\begingroup
\renewcommand{\labelenumi}{(\alph{enumi})}
\subsubsection{Economic Theory and Hypothesis}
\begin{enumerate}
    \item State a clear economic research question that your analysis will address.
    \item Provide 2--3 paragraphs reviewing relevant economic theory and citing at least 2 academic sources.
    \item Formulate a specific, testable economic hypothesis about the relationship between your key variables.
    \item Discuss the expected signs and magnitudes of coefficients based on economic theory.
\end{enumerate}

\subsubsection{Data Description and Exploratory Analysis}
\begin{enumerate}
    \item Provide detailed description of your dataset: source, time period, sample selection, variable definitions.
    \item Create comprehensive descriptive statistics table (mean, median, standard deviation, min, max, skewness, kurtosis).
    \item Generate appropriate visualizations: histograms, box plots, correlation matrix heatmap, scatter plots of key relationships.
    \item Assess normality using Q-Q plots and density plots for continuous variables.
    \item Identify and discuss potential non-linearities. What transformations might be needed?
    \item Comment on outliers using appropriate detection methods (IQR, standardized residuals).
    \item Handle missing values: document patterns, justify imputation method, sensitivity analysis.
    \item Discuss findings from Pandas Profiling analysis.
\end{enumerate}

\subsubsection{Variable Selection and Model Development}
\begin{enumerate}
    \item Apply Boruta Algorithm to identify important predictors (document parameters used).
    \item Apply Mallows $C_p$ criterion for variable selection.
    \item Check for multicollinearity using Variance Inflation Factors.
    \item Based on economic theory and statistical evidence, develop and justify three competing models:
    \begin{itemize}
        \item \textbf{Model 1:} Simple regression with your theoretically most important predictor.
        \item \textbf{Model 2:} Multiple regression with 2--3 key predictors.
        \item \textbf{Model 3:} Multiple regression including all statistically significant predictors.
    \end{itemize}
    \item Compare models using AIC, BIC, adjusted $R^2$, and cross-validation performance.
    \item Select final model and provide detailed economic and statistical justification.
\end{enumerate}

\subsubsection{Final Model Diagnostics and Robustness}
Required diagnostic tests and visualizations:
\begin{itemize}
    \item \textbf{Linearity:} Residuals vs.\ fitted values plots, RESET test for specification.
    \item \textbf{Homoskedasticity:} Breusch--Pagan test, White test, scale-location plots.
    \item \textbf{Normality of residuals:} Q-Q plots, Jarque--Bera test, histograms of residuals.
    \item \textbf{Outliers and influence:} Cook's distance plot, leverage vs.\ residuals, DFBETAS.
    \item \textbf{Robustness analysis:} Bootstrap coefficient estimates (1000 replications), compare with OLS estimates.
    \item \textbf{Model validation:} 5-fold or 10-fold cross-validation, out-of-sample prediction accuracy.
\end{itemize}
Document any violations of OLS assumptions and discuss implications for your results.

\subsubsection{Results and Economic Interpretation}
\begin{enumerate}
    \item Present final model results in professional regression table.
    \item Interpret ALL coefficients in economic terms (magnitude, direction, units).
    \item Distinguish between statistical significance and economic significance.
    \item Calculate and interpret confidence intervals for key parameters.
    \item Discuss $R^2$ and model fit in economic context.
\end{enumerate}

\subsubsection{Limitations, Policy Implications, and Extensions}
\begin{enumerate}
    \item Discuss key limitations of your analysis (data, methodology, scope).
    \item Identify potential sources of bias and how they might affect conclusions.
    \item Describe policy implications of your findings (if applicable).
    \item Suggest concrete extensions for future research.
    \item Comment on external validity: to what populations/contexts do results generalize?
\end{enumerate}

\subsubsection{Code Quality and Reproducibility}
\begin{itemize}
    \item Well-commented Python code with clear section headers.
    \item Reproducible analysis (set random seeds where applicable).
    \item Efficient, readable code following best practices.
    \item All figures and tables generated programmatically.
\end{itemize}

\endgroup
